\documentclass[a4paper,10pt]{article}
\usepackage[utf8]{inputenc}
\usepackage{graphicx}
\usepackage{subfigure}
\usepackage{fancyhdr}
\usepackage{listings}
\usepackage{ gensymb }


%opening

\pagestyle{fancy} %eigener Seitenstil
\fancyhf{} %alle Kopf- und Fußzeilenfelder bereinigen
\fancyhead[L]{SEP Abgabegespräch SS2015/16} %Kopfzeile links
\fancyhead[C]{} %zentrierte Kopfzeile
\fancyhead[R]{\includegraphics[width=1cm]{../../template/tugraz_bw}} %Kopfzeile rechts
\renewcommand{\headrulewidth}{0.4pt} %obere Trennlinie
\fancyfoot[C]{} %Seitennummer
\renewcommand{\footrulewidth}{0.4pt} %untere Trennlinie


\lstset{language=C++, basicstyle=\ttfamily\tiny}  

\begin{document}

\section*{Serienimpedanz berechnen}

Schreibe ein einfaches Programm, welches elektrische Impedanzen aus einem File ausliest und die gesamte Serienimpedanz berechnet. Eine Impedanz setzt sich aus dem Realteil $a$ und dem Imaginärteil $b$ zusammen und hat folgende Form: $\underline{Z}=a+jb$ \newline
Folgende Impedanzen stehen zur Auswahl:
\begin{description}
  \item[Widerstand R:] $\underline{Z}=a+j0$ $\rightarrow$ der Imaginärteil ist 0
  \item[Kapazität C:] $\underline{Z}=a-jb$ $\rightarrow$ der Imaginärteil ist negativ
  \item[Induktivität L:] $\underline{Z}=a+jb$ $\rightarrow$ der Imaginärteil ist positiv
\end{description}
Das File hat das Format \texttt{<typ> <Betrag von a> <Betrag von b>}.
Beispiel für eine Impedanz-Liste \texttt{impedance.txt}:
\begin{lstlisting}[frame=single]
R 10 0
L 1 2
C 5 9
R 1 10
\end{lstlisting}
Nun soll das File vom Programm eingelesen werden und der Imaginärteil entsprechend den Regeln für die Impedanzen (Widerstand R, Kapazität C, Induktivität L) gesetzt werden. Folgende Befehle stehen dem Benutzer zur Verfügung:
\begin{description}
 \item[quit] Beendet das Programm.
 \item[calculate series] Berechnet die Serienimpedanz aus.
\end{description}
Es gilt folgende Rechenvorschrift für die Serienimpedanz: $\underline{Z}=\sum\limits_{i=1}^n a_i + j\sum\limits_{i=1}^n b_i$
Mit den Werten aus der Beispieldatei ergibt sich somit: \newline$\underline{Z}=(10+1+5+1)+j(0+2-9+0)=17+j(-7)$


\subsection*{Beispiel}
\footnotesize
\begin{verbatim}
./impedance impedance.txt
> calculate series
Z=17+j-7
> quit
\end{verbatim}

\subsection*{Bewertung}
\begin{itemize}
 \item Dokumentation und Programmierstil: 1 Punkt
 \item Einhalten der OOP Konzepte: 1 Punkt
 \item Struktur \& Korrektheit: 1 Punkt
\end{itemize}
\newpage


\end{document}
